%-----------------------------------
\section{Outros tipos de trabalho}
\label{sec_tipo_outros}

Além dos tipos apresentados, outros formatos de trabalho são comuns em Sistemas de Informação. Em geral, eles combinam elementos dos tipos anteriores, e a estrutura pode ser adaptada a partir do tipo mais próximo, indicado entre parênteses.

\begin{alineas}
	\item \textbf{Design Science Research} (desenvolvimento de software e pesquisa experimental): concebe um artefato, como um modelo, método ou sistema, e o avalia com rigor quanto à sua utilidade para resolver um problema relevante \cite{hevner2004design};
	\item \textbf{pesquisa-ação} (estudo de caso): o pesquisador intervém em uma organização em ciclos de planejamento, ação, observação e reflexão, em conjunto com os participantes \cite{thiollent2011pesquisa};
	\item \textbf{modelagem e melhoria de processos de negócio} (estudo de caso): mapeia processos com notações como BPMN, identifica problemas e propõe melhorias \cite{dumas2018fundamentals};
	\item \textbf{análise de dados e \textit{Business Intelligence}} (desenvolvimento de software ou pesquisa experimental): extrai indicadores e conhecimento de dados reais, com painéis ou modelos preditivos;
	\item \textbf{avaliação de usabilidade ou acessibilidade} (pesquisa experimental ou \textit{survey}): avalia sistemas existentes com usuários ou com base em diretrizes, como as recomendações de acessibilidade para a web;
	\item \textbf{avaliação de segurança ou conformidade} (estudo de caso): analisa a aderência de uma organização ou sistema a normas e leis, como a LGPD \cite{brasil2018lgpd}. Testes de invasão só podem ser feitos com autorização formal;
	\item \textbf{plano de negócio de base tecnológica} (estudo de caso): propõe um produto ou serviço de tecnologia com análise de mercado e de viabilidade, quando aceito pelo regulamento do curso;
	\item \textbf{pesquisa qualitativa com entrevistas} (\textit{survey}): investiga percepções e experiências em profundidade, com análise de conteúdo ou temática.
\end{alineas}

Quando o trabalho não se encaixar claramente em um tipo, defina com o orientador qual estrutura melhor evidencia o problema, o método e os resultados.
